\documentclass[11pt]{article}
\usepackage[margin=1in]{geometry}
\usepackage{amsmath,amssymb,amsthm}
\usepackage{booktabs,tabularx,array}
\usepackage[T1]{fontenc}
\usepackage{lmodern}
\usepackage{enumitem}
\usepackage{xurl}
\usepackage[hidelinks]{hyperref}
\setlist{nosep,leftmargin=*}
\newcommand{\Rea}{\operatorname{Re}}

\title{The constant in [OpenAI, \emph{Uniform exclusion of Landau--Siegel zeros}] is effective: $c \ge 3.39\times 10^{-3}$}
\author{Joseph M. Reilly}
\date{Draft of 2026-10-20, not reviewed by a specialist}

\begin{document}
\maketitle

\noindent Computation: \url{src/effectivity.py}; output: \url{results/2026-10-20-effectivity/effectivity-output-v2.txt}.
Derivation: \url{docs/effectivity/derivation.md}. Sources and quotations: \url{docs/effectivity/sources.md} (S1--S8).
This PDF is built from \url{docs/effectivity/note.tex}, which follows \url{note.md}.

\section{Statement}

\textbf{Assuming the proof in [S1] is correct} (see \S\ref{sec:assumptions}), every real zero $\beta\in(0,1)$ of every primitive nonprincipal real Dirichlet $L$-function of conductor $q\ge 3$ satisfies
\[
(1-\beta)\log q > 3.399\times 10^{-3}\quad\text{using Platt's computation for } q\le 4\cdot 10^5,
\]
\[
(1-\beta)\log q > 4.467\times 10^{-3}\quad\text{using Lu--Zaman--Zhao's computation for } q\le 10^{10} \text{ instead.}
\]
The argument of [S1] alone, without any computation for small $q$, gives $(1-\beta)\log q > 5.04\times 10^{-4}$ for every $q\ge 3$ except $q=8$.
These bounds are certified in Arb ball arithmetic from exact rational parameters.

\section{What the paper proves and what is added here}

\begin{itemize}
\item \textbf{The paper (Theorem 1):} some absolute $c>0$ exists. Its proof argues by contradiction along a sequence with $q\to\infty$ and $\delta := (1-\beta)\log q\to 0$, and it leaves every analytic constant implicit ($\ll$, $C$, $C_H$).
\item \textbf{Added here:}
  \begin{itemize}
  \item every implicit constant replaced by an explicit one, with a cited source and its range of validity;
  \item the limiting argument replaced by a direct inequality valid for each $q\ge q_0$;
  \item the parameters $(\gamma,H)$ chosen to maximize the resulting bound;
  \item a verification of every side condition on which the paper's lemmas depend.
  \end{itemize}
\end{itemize}
No step of the paper's algebra (Lemmas 3, 6 and 7, Corollary 4) is changed.

\section{Explicit inputs}

{\small
\begin{tabularx}{\textwidth}{@{}>{\raggedright}p{2.6cm}>{\raggedright}X>{\raggedright\arraybackslash}p{4.4cm}@{}}
\toprule
Input & Statement used & Source and range\\
\midrule
Hadamard identity (2.3) & $\Rea(-L'/L)(s,\chi) = \tfrac12\log(q/\pi) + \tfrac12\psi((s+\varepsilon)/2) - \sum_\rho \Rea\frac{1}{s-\rho}$, terms $\ge 0$, real $s>1$ & as stated in [S1] \S2; checked against Davenport, 2nd ed., Ch.~12 (17)--(18), p.~83\\
$\zeta$ log-derivative & $-\zeta'/\zeta(\sigma) < 1/(\sigma-1)$, $\sigma>1$ & [S6] Lemma 12; its final elementary inequality proved here for all $\sigma>1$ (\S\ref{sec:direct})\\
Digamma & $\psi((s+\varepsilon)/2)\le\psi(3/2) = 2-\gamma-2\log 2$ for $1<s\le 2$ & $\psi$ increasing\\
Mertens & $\log x + E - \frac{1}{2\log x} < \sum_{p\le x}\frac{\log p}{p}$ ($x>1$); $\sum_{p\le x}\frac{\log p}{p} < \log x + E + \frac{1}{2\log x}$ ($x\ge 319$) & Rosser--Schoenfeld Thm 6, (3.21)--(3.22), p.~70, read in the original\\
Chebyshev & $\theta(x) < (1+7.5\cdot 10^{-7})x$, $x>0$ & [S5] Platt--Trudgian Cor.~2\\
Primes dividing $q$ & $\sum_{p\mid q}\frac{\log p}{p} \le \min\bigl(\tfrac12, \tfrac{\log 2}{2\log q}+\tfrac13\bigr)\log q$ & elementary\\
Lemma 6 & $S_2\le 192\,M H^{-1/3}U$, $S_1\ge M H^{2/3}U/(4\cdot 97^2)$, given $N\ge H$ and $H^{2/3}N^{4/3}\ge 2\cdot 97^2$ & [S1] lines 517--550\\
Small $q$ & no real zero with $\sigma>0$ for $q\le 4\cdot10^5$ (Platt 2016); $\delta>1/5$ for $q\le 10^{10}$ & Platt 2016 Thms 10.1--10.2, read in the original; [S2] Theorem 1.1\\
\bottomrule
\end{tabularx}
}

\section{The direct inequality and the parameter choice}\label{sec:direct}

Combining the explicit Lemma 2, (5.1), (5.6) with the $\theta$ bound, and Lemma 6 (derivation.md \S\S1--5): if a real zero with $\delta=(1-\beta)\log q$ exists and $N\ge H$, then
\[
1 \le \tfrac34(1+\rho_H) + \frac{C_\ell+\frac12(1+\rho_H)}{\lambda} + \frac e2\,\delta\lambda + \frac{C_H+(1+\rho_H)\log 8}{L} + \frac{(1+7.5\cdot10^{-7})\cdot 4\cdot 97^2}{H^{2/3}L} + \frac{6\cdot 97^2}{H^{2/3}U},
\]
where
\begin{itemize}
\item $\rho_H = 7{,}226{,}112/H$, $C_\ell = e/4 + c_q$, $C_H = \log H + \frac{1}{2\log H} + \frac{1}{2L} + \frac e2\kappa$;
\item $\kappa = \tfrac12\psi(3/2) - \tfrac12\log\pi = -0.55412\ldots < 0$;
\item $L = \log U = \tfrac43\log N$ and $\lambda = L/\log q$.
\end{itemize}
Take $N=\lceil q^\gamma\rceil$. For all $q\ge q_0$:
\begin{itemize}
\item $\lambda$ lies in $[4\gamma/3,\ 4\gamma/3 + \tfrac43\log(1+q_0^{-\gamma})/\log q_0]$;
\item every term except the $\delta$ term is largest at $q=q_0$;
\item so every real zero with $q\ge q_0$ has $\delta>\delta_0(q_0,\gamma,H)$, with $\delta_0$ given in closed form (function \texttt{delta0} in \url{src/effectivity.py}).
\end{itemize}
The paper's split of the margin ($13/16+1/16$) is not needed: for fixed $(\gamma,H)$, $\delta_0$ is determined exactly.

The parameters were found by a floating-point search and then certified in Arb at the exact values below (\url{effectivity-output-v2.txt}; every printed bound is the Arb lower endpoint rounded down).

\begin{center}
\begin{tabular}{@{}lrrrlr@{}}
\toprule
Regime & $q_0$ & $\gamma$ & $H$ & $c_q$ & certified $\delta_0 >$\\
\midrule
Platt & 400{,}001 & 201/10 & 1{,}995{,}262{,}314 & $\frac{\log 2}{2\ell}+\frac13$ & $\mathbf{3.3992904026\times10^{-3}}$\\
Lu--Zaman--Zhao & $10^{10}+1$ & 153/10 & 2{,}818{,}382{,}931 & $\frac{\log 2}{2\ell}+\frac13$ & $\mathbf{4.4673780338\times10^{-3}}$\\
argument alone & 3 & 671/5 & 1{,}122{,}018{,}454 & $\frac12$ & $\mathbf{5.0425352961\times10^{-4}}$\\
\bottomrule
\end{tabular}
\end{center}

At the paper's own $H=86{,}713{,}344$, the same three regimes give 2.064, 2.652 and 0.3245 ($\times10^{-3}$). Optimizing over $H$ is worth a factor of about 1.6.

\paragraph{Side conditions} (derivation.md \S7), certified at $q_0$ and implied for all $q\ge q_0$:
$H\ge 319$; $N\ge H$ ($\gamma\log q_0>\log H$); $H^{2/3}N^{4/3}\ge 2\cdot97^2$ (from $N\ge H\ge 319$); $U\ge 3$, and $s=1+1/L\in(1,2]$; $\kappa<0$.

\paragraph{Lemma 12's last step.}
$\varphi(\sigma) = \frac{1}{2\sigma-1} - \Bigl[\frac{1}{\sigma-1} - \frac{1-2^{-\sigma}}{\sigma} - \frac{3^{1-\sigma}}{\sigma-1}\Bigr] > 0$ for all $\sigma>1$. Proved in three pieces:
\begin{itemize}
\item on $(1,1.01]$, by the exact chain $(1-3^{-t})/t\le\log 3$ and $(1-2^{-\sigma})/\sigma\ge 1/(2\sigma)$, which reduces it to $\log 3 - 1/2.02 < 1/1.02$;
\item on $[1.01,5]$, by Arb on 4{,}000 intervals (minimum certified lower bound of $\varphi$: $0.03878$);
\item for $\sigma\ge 5$, by $\varphi>0 \Leftarrow \sigma^2-5\sigma+2>0$.
\end{itemize}

\section{Small conductors}

\begin{itemize}
\item \textbf{Platt regime:} for $q\le 4\cdot10^5$ there is no real zero with $\sigma>0$ at all, so the bound holds vacuously. This range includes $q=8$, the conductor the paper's field construction excludes ($d=2$). $q>4\cdot10^5$ is covered by the first row.
\item \textbf{Lu--Zaman--Zhao regime:} for $q\le10^{10}$, $\delta>1/5$, which also covers $q=8$; $q>10^{10}$ is covered by the second row; so $c=\min(4.467\cdot10^{-3},1/5)=4.467\cdot10^{-3}$.
\item \textbf{Argument alone:} covers $q\ge3$ with $q\ne8$ by taking $\gamma$ large ($\lambda\approx179$), i.e.\ $N\approx q^{134}$. $q=8$ needs a separate input.
\end{itemize}

\section{The final constant}

\[
c = 3.399\times10^{-3}\ \text{(with Platt)},\qquad c = 4.467\times10^{-3}\ \text{(with Lu--Zaman--Zhao)},
\]
both rounded down from the certified values.

\paragraph{Context.}
\begin{itemize}
\item The companion paper on the zero-free half-plane $\Rea s>7/8$ would, if correct, give $1-\beta\ge 1/8$ for every such $\beta$, a far stronger statement.
\item The point here is narrower: the nine-page Landau--Siegel argument, on its own, yields an explicit constant.
\item A real zero with $\delta\approx0.003$ is far outside anything ever observed. These numbers are for checking the argument, not a claim about zeros.
\end{itemize}

\paragraph{Comparison with an earlier version} (docs/explicit-constant.md, 2026-10-08), whose figures were $3.0379\cdot10^{-4}$, $1.99806\cdot10^{-3}$ and $2.58930\cdot10^{-3}$:
\begin{itemize}
\item \textbf{Same inputs reproduce it.} With its inputs switched back on, this code gives slightly larger maxima ($3.040\cdot10^{-4}$, $2.0000\cdot10^{-3}$, $2.5918\cdot10^{-3}$, rounded down). The earlier note reported a $\delta$ with a positive safety margin ($F\approx9.4\cdot10^{-5}$) rather than the maximum.
\item \textbf{Its ``all $q\ge3$'' figure} covered small $q$ the same way as row 3 here: by taking $\lambda$ large, not by any computation ($q=8$ again excepted).
\item \textbf{It used three inputs replaced here:} an unsourced ``$\sum_{p\le x}\log p/p<\log x$'' for the primes up to $H$, which costs $|E|=1.33$ in $C_H$ against the sourced Rosser--Schoenfeld upper bound; Rosser--Schoenfeld Thm 9 for $\theta$; a fixed $H$.
\item \textbf{The gain comes mainly from freeing $H$} ($\times1.6$). Of the other changes, $c_q=\frac12$ is worse than its $c_q$ at large $q$ but better at $q=3$, and the sourced Mertens bound and the $\theta$ bound are improvements.
\end{itemize}

\section{Assumptions (everything the constant depends on)}\label{sec:assumptions}

\begin{enumerate}
\item \textbf{The paper's proof is correct:} in particular Lemma 3, Corollary 4, Lemma 6, Lemma 7 and the bounds (5.1) and (5.6).
  \begin{itemize}
  \item This note does not re-verify them mathematically.
  \item Independent evidence: the Lean formalization of Theorem 1 is accepted by the official comparator with two kernels (Lean and nanoda), including a version rerouted through the paper's own Lemma 3. See the replay packet (DOI 10.5281/zenodo.23269527) and the statement audit on \url{grwtsk/openai-math} \#40.
  \item The Lean proof has existential constants only and does not contain these numbers.
  \end{itemize}
\item \textbf{The Hadamard identity (2.3).} Checked against Davenport, \emph{Multiplicative Number Theory}, 2nd ed., Ch.~12, (17) and (18) on p.~83, with $a=\varepsilon$ (p.~70). Taking real parts of (17) and using (18), $\Rea B(\chi) = -\sum\Rea 1/\rho$, gives (2.3).
\item \textbf{Rosser--Schoenfeld Theorem 6,} read in the original: (3.21) for $1<x$ and (3.22) for $319\le x$, p.~70. $E\approx-1.33258\,22757\,33221$, (2.11). This agrees with the quotation in arXiv:1703.08032 and with the enclosure in \url{src/effectivity.py}.
\item \textbf{Platt (2016),} read in the original. Theorem 10.1 states GRH for primitive characters of modulus $q\le 400\,000$, up to a positive height. Theorem 10.2 states $L(1/2,\chi)\ne0$ for $q\le 2\,000\,000$. Together (our inference) they rule out a real zero in $(0,1)$ for $q\le4\cdot10^5$, matching the Lu--Zaman--Zhao summary. Needed only for the $3.399\cdot10^{-3}$ statement.
\item \textbf{Lu--Zaman--Zhao Theorem 1.1} (arXiv:2602.03626v1, a preprint), read in the original. Needed only for the $4.467\cdot10^{-3}$ statement.
\item \textbf{Platt--Trudgian Corollary 2,} read in the original.
\item \textbf{Lemma 12 of arXiv:2511.13898} (a preprint), read in the original. Its reduction steps are taken from that proof; its final inequality is proved here.
\item \textbf{$\chi(p)=(d/p)$ for odd $p\nmid q$, and $|d|\le q$,} as stated in the paper.
\end{enumerate}
Not used: any unproved hypothesis (GRH or similar), or any floating-point value in a claimed inequality.

\section*{References}
\begin{enumerate}[label={[S\arabic*]}]
\item OpenAI, \emph{Uniform exclusion of Landau--Siegel zeros}, October 1, 2026, \url{github.com/openai/math}, commit adc7f12, family 003.
\item R.~F.~Lu, A.~Zaman, H.~Zhao, \emph{Numerical computations concerning Landau--Siegel zeros}, arXiv:2602.03626v1 (2026).
\item J.~B.~Rosser, L.~Schoenfeld, \emph{Approximate formulas for some functions of prime numbers}, Illinois J.\ Math.\ 6 (1962), 64--94, \url{doi:10.1215/ijm/1255631807}.
\item Axler, \emph{New estimates for some functions defined over primes}, arXiv:1703.08032.
\item Platt, Trudgian, \emph{On the first sign change of $\theta(x)-x$}, arXiv:1407.1914.
\item Upper bounds on gaps between zeros of $L$-functions, arXiv:2511.13898, Appendix, Lemma 12.
\item H.~Davenport, \emph{Multiplicative Number Theory}, 2nd ed., rev.\ H.~L.~Montgomery, Springer GTM 74.
\end{enumerate}
\noindent Also: D.~J.~Platt, \emph{Numerical computations concerning the GRH}, Math.\ Comp.\ 85 (2016), 3009--3027, doi:10.1090/mcom/3077.

\end{document}
